\thispagestyle{empty}
\begin{tcolorbox}[enhanced, colback=hblue, colframe=hblue, boxrule=0pt, arc=4pt, left=16pt, right=16pt, top=22pt, bottom=22pt]
  \color{white}\sffamily
  {\large\bfseries Bewijzen in de Wiskunde \textperiodcentered\ 2026--2027}\\[10pt]
  {\fontsize{30}{34}\selectfont\bfseries Planning exercises}\\[4pt]
  {\fontsize{20}{24}\selectfont worked step by step \textperiodcentered\ Chapters 0--5}\\[14pt]
  {\normalsize Every essential (\fE) and recommended (\fR) exercise of the course planning for Chapters 0--5, written out in full: what is given, what has to be shown, the plan, every step with the rule or definition that justifies it, and a checklist to compare with your own attempt.}
\end{tcolorbox}

\vspace{10pt}
\begin{tcolorbox}[enhanced, colback=white, colframe=hblue2, boxrule=0.8pt, arc=2pt, left=10pt, right=10pt, top=8pt, bottom=8pt]
\sffamily\small\sloppy
\begin{itemize}[leftmargin=1.2em]
\item \textbf{\textcolor{hblue}{Source.}} C.~Newstead, \emph{An Infinite Descent into Pure Mathematics}, Utrecht adaptation (29 July 2026). All numbers (Definition 0.36, Strategy 1.1.28, Theorem 0.42, \dots) are the book's. Nothing is used that the book has not stated before the exercise, except the elementary algebra the book takes for granted.
\item \textbf{\textcolor{hblue}{Which exercises.}} Exactly the exercises flagged \fE\ essential and \fR\ recommended in the course planning (plus the three bonus ones), in book order. The companion \emph{Workbook Ch.\,0--5} contains the theory and shorter solutions of \emph{all} exercises.
\item \textbf{\textcolor{hblue}{How to use it.}} Do the exercise yourself first. Then read the solution top to bottom and tick the \emph{Compare with your work} list: it names the sentences a complete answer must contain and the places where marks are usually lost.
\end{itemize}
\end{tcolorbox}

\vspace{8pt}
\hsub{How each solution is laid out}

\begin{center}\sffamily\small
\begin{tabular}{@{}p{0.26\linewidth}p{0.69\linewidth}@{}}
\textcolor{hblue}{\textbf{Given.}} & Every assumption, including the silent ones: what kind of object each letter is, and what the hypotheses mean once their definitions are unfolded.\\[3pt]
\textcolor{hblue}{\textbf{To show.}} & The goal, often rewritten in symbols, with its \emph{logical shape} (an implication, a ``for all'', an ``if and only if'', a negation, \dots). The shape decides the first sentence of the proof.\\[3pt]
\textcolor{hgreenline}{\textbf{Plan.}} & The strategy from the book that fits the shape, and the one idea that makes the proof work.\\[3pt]
\textcolor{hblue}{\textbf{Step 1, Step 2, \dots}} & The proof, one move per step. Grey text in brackets, like \why{Definition 0.36}, names the definition, strategy, theorem or algebra rule that justifies the move. Variables are introduced the moment they appear.\\[3pt]
\textcolor{hblue}{\textbf{Conclusion.}} & The last sentence, closing every open ``suppose'' and ``let'', followed by $\square$.\\[3pt]
\textcolor{hpurple}{\textbf{Compare with your work.}} & What a complete answer must contain, and the typical mistake for this exercise.
\end{tabular}
\end{center}

\begin{key}[{Conventions of the book, used throughout}]
\begin{itemize}
\item $\N=\{0,1,2,\dots\}$ (Definition 0.6); $[n]=\{k\in\N\mid k<n\}$ (Definition 0.25); intervals consist of real numbers (Definition 0.28).
\item ``$b$ divides $a$'', written $b\mid a$, means $a=qb$ for some $q\in\Z$ (Definition 0.36). It is a statement about multiplication: never write the fraction $a/b$ in a divisibility proof.
\item ``Even'' means divisible by $2$; ``odd'' means not even (Definition 0.41). That an odd $a$ equals $2b+1$ is Proposition 1.1.58, to be quoted.
\item Division theorem (Theorem 0.42): for $b\ne0$ there is \emph{exactly one} pair $q,r\in\Z$ with $a=qb+r$ and $0\le r<|b|$. The word ``exactly one'' is what lets you identify a remainder.
\item Truth tables use $\checkmark$ and $\times$; $\bot$ is a contradiction; $\equiv$ is logical equivalence (a statement \emph{about} formulae).
\item In every proof the \emph{reader} chooses the arbitrary element (``let $x\in X$''); you choose the witness of an existential (``define $n=\dots$'').
\end{itemize}
\end{key}

\begin{key}[{The flags}]
\fE\ essential in the planning \quad \fR\ recommended \quad bonus = marked bonus in the planning. Quiz 1 (23 Sep): numbers and \S1.1. Quiz 2 (7 Oct): sets and functions. Quiz 3 (21 Oct): relations and modular arithmetic. Exam 4 Nov.
\end{key}
